\subsection{齐次主集}

定义 7. 设 G 为一个群。G 在集合 E 上的一个作用称为单传递的，如果存在 E 中的一个元素 x，使得由 x 定义的轨道映射是双射。集合 E 连同 G 在 E 上的一个单传递左作用，称为 G 下的左齐次主 G-集（或 G 下的左齐次主集）。

这等同于说 G 在 E 上自由且传递地作用，或者也说存在一个元素 \(x \in E\) 使得由 x 定义的轨道映射是 G-集 G（其中 G 通过左平移作用）到 E 的一个同构；或者等价地，满足以下两个条件：

\begin{enumerate}
\def\labelenumi{(\roman{enumi})}
\item
  E 是非空的。
\item
  对于所有元素 \(x\) 和 \(y\) 于 \(\mathbf{E}\) ，存在且仅存在一个元素 \(\alpha \in \mathbf{G}\) 的集合，使得 \(\alpha x = y\) .
\end{enumerate}

条件(ii)也等价于以下条件：

\begin{enumerate}
\def\labelenumi{(\roman{enumi})}
\setcounter{enumi}{2}
\tightlist
\item
  映射 \((\alpha, x) \mapsto (\alpha x, x)\) 是从 \(G \times E\) 到 \(E \times E\) .
\end{enumerate}

我们将定义右齐次主 G-集的任务留给读者。

例子。(1) 设 G 通过左（相应地，右）平移作用于自身。这样便在 G 上定义了一个左（相应地，右）齐次主 G-集结构，它有时记为 \(G_{s}\) （相应地，\(G_{d}\) ）。

\begin{enumerate}
\def\labelenumi{(\arabic{enumi})}
\setcounter{enumi}{1}
\item
  设 E 是交换群 G 下的齐次集。若 G 忠实作用于 E，则后者是齐次主 G-集。
\item
  设 E 和 F 是两个同构的、带同一种类结构的集合，令 \(\operatorname{Isom}(\mathbf{E}, \mathbf{F})\) 为 E 到 F 上（带给定结构的）同构所成的集合。E 的（带给定结构的）自同构群 \(\operatorname{Aut}(\mathbf{E})\) 在 \(\operatorname{Isom}(\mathbf{E}, \mathbf{F})\) 上依法则 \((\sigma, f) \mapsto f \circ \sigma\) 右作用，且 \(\operatorname{Isom}(\mathbf{E}, \mathbf{F})\) 是一个右齐次主 \(\operatorname{Aut}(\mathbf{E})\) -集。类似地，群 \(\operatorname{Aut}(\mathbf{F})\) 在 \(\operatorname{Isom}(\mathbf{E}, \mathbf{F})\) 上依法则 \((\sigma, f) \mapsto \sigma \circ f\) 左作用，且 \(\operatorname{Isom}(\mathbf{E}, \mathbf{F})\) 是一个左齐次主 \(\operatorname{Aut}(\mathbf{F})\) -集。
\item
  *向量空间的加法群下的齐次主集称为仿射空间（参见 II，§ 9，第 1 号）。*
\end{enumerate}

齐次主 G-集 \(G_{s}\) （例 1）的自同构群是 G 的右平移群，它等同于 \(G^{0}\) （第 5 号，命题 5）。设 E 是一个齐次主 G-集，a 是 E 的一个元素。由 a 定义的轨道映射 \(\omega_{a}\) 是 G-集 \(G_{s}\) 到 E 上的同构。通过结构转移，得到一个同构 \(\psi_{a}\) ，它把 \(G^{0}\) 映到 Aut(E) 上。应当注意， \(\psi_{a}\) 一般依赖于 a；更确切地说，对于 \(\alpha \in G\) ，

\[
\psi_ {\alpha a} = \psi_ {a} \circ \operatorname{Int} _ {\mathbf {G} ^ {0}} (\alpha) = \psi_ {a} \circ \operatorname{Int} (\alpha^ {- 1}).\tag{3}
\]

因为，用 \(\delta_{a}\) 表示 G 上的平移 \(x\mapsto x\alpha\) ，

\[
\omega_ {\alpha a} = \omega_ {a} \circ \delta_ {\alpha}
\]

和

\[
\psi_ {a} (x) = \omega_ {a} \circ \delta_ {x} \circ \omega_ {a} ^ {- 1}, \quad x \in G,
\]

因此

\[
\psi_ {\alpha a} (x) = \omega_ {a} \circ \delta_ {\alpha} \circ \delta_ {x} \circ \delta_ {\alpha} ^ {- 1} \circ \omega_ {a} ^ {- 1} = \omega_ {a} \circ \delta_ {\alpha^ {- 1} x \alpha} \circ \omega_ {a} ^ {- 1} = \psi_ {a} (\alpha^ {- 1} x \alpha).
\]

